% page 594 du fac-simile (image p-593 du scan)
% bloc 157.5 x 233.3 mm ; marge gauche 8.0 mm
\pgc{-12.700mm}{0.000mm}{
\l{\cel{3}586.}
\l{}
\l{\sou{torida}. Ube atmosfero esas extreme varma. La\cel{1}\sou{zono torida}:}
\l{\cel{3}la parto di la Terglobo, inter la du tropiki. - EFIS.}
\l{}
\l{\sou{tormentar}. (\sou{trans.})\cel{2}I. Submisar (kondamnito) a korpo-puniso}
\l{\cel{3}tre granda, grava, kruela, netolerebla.- II. Probar}
\l{\cel{3}arachar de pre-akuzato konfesi, per tordar, luxacar, dis-}
\l{\cel{3}lokar, brular, kombustar, ca o ta parto di lua korpo.}
\l{\cel{3}- EFIS.}
\l{}
\l{\sou{tormentilo}. (\sou{bot.)}\cel{2}Planto ek la familio \textquotedbl{}rozacei\textquotedbl{}, di qua}
\l{\cel{3}la radiko esas astriktiva. - L.\cel{1}\sou{tormentilla}.}
\l{}
\l{\sou{tornar}. (\sou{trans.}) Fasonar per mashino sur qua onu muntis (ver-}
\l{\cel{3}tikale od horizontale) peco de ligno, de metalo, de ivoro,}
\l{\cel{3}e c., quan onu rotacigas, tale ke rivenas, pos intertempi}
\l{\cel{3}inter-egala, ta o ca parto di la peco por subisar la efiko}
\l{\cel{3}da cizelo od irg-altra-sorta utensilo konceptita por faso-}
\l{\cel{3}nar lu. - FIS.}
\l{}
\l{\sou{tornasolo}. I. Materio kolorizanta, blua, quan onu extraktas}
\l{\cel{3}de varietato di krotono (sunfloro, L.\cel{1}\sou{rocella lecanora}.}
\l{\cel{3}- II. Solvuro de ta blua materio, quan la acidi igas red-}
\l{\cel{3}eskar. - EFS.}
\l{}
\l{\sou{torono}. (\sou{teknol.}) Asemblajo ek fili kuntordita por facar kor-}
\l{\cel{3}degi, kabli, e c., e di qua la nombro varias segun la gro-}
\l{\cel{3}seso di la kordego, di la kablo. - FI.}
\l{}
\l{\sou{torpedo.}\cel{1}I. (\sou{zool.}) Fisho kartilagoza, kun korpo plata, disko-}
\l{\cel{3}forma, qua havas, an singla latero di sua kapo, organo}
\l{\cel{3}partikulara qua elektro-sukusas ti qui tushas ta fisho. -}
\l{\cel{3}II. (\sou{milit-arto}). Mashino plenigita ye pulvero, kun amorc-}
\l{\cel{3}ilo, qua, se shokata, explozas, e quan onu lansas kontre}
\l{\cel{3}la navi enemikala, od imersas, por impedar la proximesko}
\l{\cel{3}da ta navi a la litoro, e c. - DEFIRS.}
\l{}
\l{\sou{torporar}. (\sou{netrans}.) Afektesar da ta sorto di inerteso, di}
\l{\cel{3}graveso quan produktas la ceso di la sentiveso e di la}
\l{\cel{3}funciono, en un o plura parti di la homo-korpo, o di la}
\l{\cel{3}psiko. - EFIS.}
\l{}
\l{\sou{torso}. (\sou{anat}.)Korpo di homo o di animalo, sen la kapo nek la}
\l{\cel{3}membri. - DEFIR.}
\l{}
\l{\sou{torsado}. I. Volvajo ek fili de silko, oro, e c., tordita spi-}
\l{\cel{3}rale, de qui onu facas la kordeti di la franji, di la}
\l{\cel{3}glani por kurteni, tapeti, e c. - II. Volvajo ek fili de}
\l{\cel{3}arjento, de oro, de qua onu facas la epoleti di la oficiri.}
\l{\cel{3}- III. Volvajo de hari tordita spirale. - DF.}
\l{}
\l{\sou{torto.}\cel{2}Kukajo ronda, en qua onu pozas karno de bruti, de}
\l{\cel{3}pultro, de fisho - e kaneli. - DFIRS.}
}
